\begin{enumerate}
  \item Les respostes s'han d'ajustar a l'enunciat de la pregunta. Es valorarà sobretot que l'alumnat demostri que té clars els conceptes de caràcter físic sobre els quals tracta cada pregunta.
  \item Es tindrà en compte la claredat en l'exposició dels conceptes, dels processos, dels passos a seguir, de les hipòtesis, l'ordre lògic, l'ús correcte dels termes científics i la contextualització segons l'enunciat.
  \item En les respostes cal que l'alumnat mostri una adequada capacitat de comprensió de les qüestions plantejades i organitzi de forma lògica la resposta, tot analitzant i utilitzant les variables en joc. També es valorarà el grau de pertinença de la resposta, el que l'alumnat diu i les mancances manifestes sobre el tema en qüestió.
  \item Totes les respostes s'han de raonar i justificar. Un resultat erroni amb un raonament correcte es valorarà. Una resposta correcta sense raonament ni justificació pot ser valorada amb un 0, si el corrector no és capaç de veure d'on ha sortit el resultat.
  \item Tingueu en compte que un error no s'ha de penalitzar dues vegades en el mateix problema. Si un apartat necessita un resultat anterior, i aquest és erroni, cal valorar la resposta independentment del seu valor numèric, i tenir en compte el procediment de resolució.
  \item Si la resolució presentada a l'examen és diferent però correcta i està d'acord amb els requeriments de l'enunciat, s'ha d'avaluar positivament encara que no coincideixi amb la resolució donada a la pauta de correcció.
  \item Un o més errors en les unitats d'un apartat restarà 0,25 punts en la qualificació d'aquest l'apartat. Es consideren errors d'unitats: ometre les unitats en els resultats (finals o intermedis), utilitzar unitats incorrectes per una magnitud (tant en els resultats com en els valors intermedis) o operar amb magnituds d'unitats incompatibles (excepte en el cas d'un quocient on numerador i denominador tenen les mateixes unitats). Exemple: si l'apartat (a) val 1,25 punts i només s'ha equivocat en les unitats l'haurem de puntuar amb 1 punt.
  \item Un o més errors de càlcul en un apartat restarà 0,25 punts en la qualificació d'aquest apartat. Exemple: si l'apartat (a) val 1,25 punts i només s'ha equivocat en els càlculs l'haurem de puntuar amb 1 punt.
  \item Cal resoldre els exercicis fins al resultat final i no es poden deixar indicades les operacions.
  \item Cal fer la substitució numèrica en les expressions que s'utilitzen per resoldre les preguntes.
  \item Un resultat amb un nombre molt elevat de xifres significatives (6 xifres significatives) es penalitzarà amb 0,1p.
\end{enumerate}
